\pagenumbering{roman}
\hypersetup{pageanchor=true}
\newcommand{\authorsignaturearea}{%
\begin{center}
\begin{tabular}{>{\centering\arraybackslash}m{0.46\textwidth}>{\centering\arraybackslash}m{0.46\textwidth}}
\textit{Tác giả 1} & \textit{Tác giả 2} \tabularnewline[0.15cm]
\textit{(Ký tên và ghi rõ họ tên)} & \textit{(Ký tên và ghi rõ họ tên)} \tabularnewline[1.45cm]
Phí Vương Tường Tâm & Vũ Hưng Thịnh \tabularnewline
\end{tabular}
\end{center}}
\chapter*{LỜI CẢM ƠN}
Nhóm xin cảm ơn giảng viên Dung Cẩm Quang đã hướng dẫn đề tài trong học phần Dự án Công nghệ thông tin. Nhóm cũng cảm ơn các tác giả đã công bố phương pháp, mã nguồn và dữ liệu phục vụ nghiên cứu trích xuất thông tin từ tài liệu.

Bản thảo tổng hợp quá trình khảo sát, triển khai và đánh giá các mô hình đến ngày 14 tháng 11 năm 2026.

\begin{flushright}\textit{TP. Hồ Chí Minh, ngày 14 tháng 11 năm 2026}\end{flushright}
\authorsignaturearea

\chapter*{CÔNG TRÌNH ĐƯỢC HOÀN THÀNH\\TẠI TRƯỜNG ĐẠI HỌC TÔN ĐỨC THẮNG}
Chúng em xin cam đoan đây là công trình của nhóm, được thực hiện dưới sự hướng dẫn của giảng viên Dung Cẩm Quang. Số liệu thực nghiệm được tổng hợp từ các lượt huấn luyện và đánh giá đã lưu. Các nội dung và hình ảnh sử dụng từ công trình khác được ghi rõ nguồn.

Chúng em chịu trách nhiệm về nội dung báo cáo và những vi phạm tác quyền, bản quyền do nhóm gây ra trong quá trình thực hiện, nếu có.

\begin{flushright}\textit{TP. Hồ Chí Minh, ngày 14 tháng 11 năm 2026}\end{flushright}
\authorsignaturearea

\chapter*{TÓM TẮT}
Đề tài so sánh sáu mô hình trích xuất thông tin từ hóa đơn gồm RoBERTa, BROS, PICK, LayoutLMv3, Donut và Qwen2.5-VL. Các mô hình có đầu vào từ văn bản đơn thuần đến ảnh hóa đơn đi kèm chỉ dẫn. Nhóm đánh giá trên SROIE, CORD-v2 và WildReceipt bằng Precision, Recall và Macro-F1 ở mức trường--giá trị.

Đã có kết quả cho 16 trong 18 cặp mô hình và bộ dữ liệu. Trong số các lượt hoàn tất, LayoutLMv3 dẫn đầu trên SROIE (0,8869), RoBERTa dẫn đầu trên CORD-v2 (0,7984), và BROS dẫn đầu trên WildReceipt (0,9094). Qwen2.5-VL đạt Macro-F1 0,8381 trên 347 mẫu test SROIE. Hai lượt Qwen còn lại đang chờ kết quả hoặc chưa chạy. Các nhận xét áp dụng cho cấu hình thực nghiệm và cách chấm được mô tả trong báo cáo.

Từ khóa -- trích xuất thông tin, hóa đơn, mô hình thị giác--ngôn ngữ, Macro-F1.

\chapter*{ABSTRACT}
This project compares six models for receipt information extraction: RoBERTa, BROS, PICK, LayoutLMv3, Donut, and Qwen2.5-VL. The models process inputs ranging from text alone to images accompanied by instructions. The evaluation uses SROIE, CORD-v2, and WildReceipt, with field-level macro precision, recall, and F1 based on exact value matches after normalization.

Results are available for 16 of the 18 model--dataset combinations. Among completed runs, LayoutLMv3 leads on SROIE (0.8869), RoBERTa on CORD-v2 (0.7984), and BROS on WildReceipt (0.9094). Qwen2.5-VL obtains a macro F1 of 0.8381 on 347 SROIE test samples. Its CORD-v2 and WildReceipt test results are pending. The comparison is specific to the evaluated configurations and scoring protocol.

Keywords -- information extraction, receipts, vision-language models, macro F1.

\clearpage
\begingroup
\renewcommand{\baselinestretch}{1.05}\normalsize
\tableofcontents
\endgroup
\clearpage
\phantomsection\addcontentsline{toc}{chapter}{\listfigurename}\listoffigures
\clearpage
\phantomsection\addcontentsline{toc}{chapter}{\listtablename}\listoftables
\chapter*{DANH MỤC CÁC CHỮ VIẾT TẮT}
\addcontentsline{toc}{chapter}{DANH MỤC CÁC CHỮ VIẾT TẮT}
\begin{table}[H]
\centering
\renewcommand{\arraystretch}{1.25}
\begin{tabularx}{\textwidth}{lX}
\toprule
\textbf{Viết tắt} & \textbf{Ý nghĩa} \\
\midrule
KIE & Key Information Extraction -- trích xuất thông tin chính \\
OCR & Optical Character Recognition -- nhận dạng chữ \\
GT & Ground truth -- giá trị tham chiếu \\
BIO & Begin, Inside, Outside -- vị trí của token trong thực thể \\
CRF & Conditional Random Field -- trường ngẫu nhiên có điều kiện \\
LoRA & Low-Rank Adaptation -- tinh chỉnh hạng thấp \\
JSON & JavaScript Object Notation -- định dạng dữ liệu có cấu trúc \\
KV cache & Bộ nhớ lưu các tensor khóa và giá trị trong quá trình sinh \\
\bottomrule
\end{tabularx}
\end{table}
\clearpage
\pagenumbering{arabic}

